\documentclass[11pt]{amsart}
\usepackage[margin=1.1in]{geometry}
\usepackage{amsmath,amssymb,amsthm}
\usepackage{booktabs}
\usepackage{alltt}
\usepackage[colorlinks=true,linkcolor=blue,citecolor=blue,urlcolor=blue]{hyperref}

\newtheorem{theorem}{Theorem}
\newtheorem{lemma}[theorem]{Lemma}
\newtheorem{proposition}[theorem]{Proposition}
\newtheorem{corollary}[theorem]{Corollary}
\theoremstyle{definition}
\newtheorem{definition}[theorem]{Definition}
\theoremstyle{remark}
\newtheorem{remark}[theorem]{Remark}

\newcommand{\C}{\mathbb{C}}
\newcommand{\R}{\mathbb{R}}
\newcommand{\Tr}{\operatorname{Tr}}
\newcommand{\Herm}{\operatorname{Herm}}
\newcommand{\SN}{\operatorname{SN}}
\newcommand{\one}{\mathbb{I}}
\newcommand{\tr}{\mathrm{tr}}
\newcommand{\ket}[1]{|#1\rangle}
\newcommand{\bra}[1]{\langle #1|}
\newcommand{\lean}[1]{\texttt{#1}}
\setlength{\emergencystretch}{3em}

\begin{document}

\title[Tavakoli--Morelli Schmidt-number conjectures]{Proof of two conjectures of Tavakoli and Morelli\\ on Schmidt-number witnesses}

\author{Ansh Mishra}
\author{Aryan Senthilkumar}

\date{September 2026}

\begin{abstract}
Tavakoli and Morelli [Phys.\ Rev.\ A \textbf{110}, 062417 (2024)] proposed Schmidt-number witnesses based on the
trace norm of the matrix of joint outcome probabilities of two local measurements. They conjectured tight bounds
for incomplete sets of mutually unbiased bases (their Conjecture~3) and for equiangular measurements (their
Conjecture~2). We prove both conjectures. For $m$ mutually unbiased bases in $\C^d$ on each side and every state
of Schmidt number at most $r$, the correlation matrix satisfies $\|Q_m\|_{\tr}\le 1+(m-1)r/d$, and this is tight
for every $d$, $m$ and $r$. For equiangular measurements with $n>d$ outcomes,
$\|P_n\|_{\tr}\le[d(d-1)+r(n-d)]/[n(n-1)]$. Both results follow from a single inequality for rank-one
measurements whose frame operator fixes the identity. Its proof combines a Gram factorisation of the correlation
matrix, H\"older's inequality for Schatten norms, and a Bessel-type bound for the measurement frame. Both theorems
are formalised in the Lean~4 proof assistant with the Mathlib library, using only Lean's standard axioms.
\end{abstract}

\maketitle

\section{Introduction}

The Schmidt number of a bipartite state is the smallest $r$ such that the state is a mixture of pure states of
Schmidt rank at most $r$~\cite{TerhalHorodecki2000}. It quantifies the dimensionality of entanglement, and
certifying it from few measurements is a basic task in high-dimensional quantum information; see, e.g.,
\cite{Spengler2012,Bavaresco2018}. Tavakoli and Morelli~\cite{TavakoliMorelli2024} proposed witnesses built from
the matrix of joint outcome probabilities when Alice and Bob each measure one of several local bases or a single
informationally rich measurement. The criterion is an upper bound on the trace norm of this matrix, valid for all
states of Schmidt number at most $r$; a larger trace norm certifies Schmidt number larger than $r$. For
incomplete sets of mutually unbiased bases (MUBs)~\cite{WoottersFields1989,Durt2010} and for equiangular
measurements they conjectured the tight value of this bound (Conjectures~3 and~2 of~\cite{TavakoliMorelli2024}).

We prove both conjectures (Theorems~\ref{thm:mub} and~\ref{thm:eam}). Both are corollaries of a general
inequality (Theorem~\ref{thm:general}): if the frame operator of a rank-one measurement maps the identity to a
multiple of itself, with eigenvalue $\alpha$, and is bounded by $\beta$ on traceless operators, then states of
Schmidt number at most $r$ give correlation matrices of trace norm at most $\beta+(\alpha-\beta)r/d$ (when both
parties use such measurements with the same constants). The proof has two ingredients. First, a Gram
factorisation of the correlation matrix of a pure state, combined with H\"older's inequality, bounds its trace
norm by a product of two ``index of coincidence'' quantities, one for each party (Lemma~\ref{lem:factor}).
Second, a Bessel-type bound controls these quantities through the trace and the Hilbert--Schmidt norm of the
modulus of the coefficient matrix (Lemmas~\ref{lem:frameMUB} and~\ref{lem:frameEAM}). The Schmidt rank enters only
through the elementary inequality $\|X\|_1^2\le r\|X\|_2^2$.

A natural first attempt applies the triangle inequality to the Schmidt decomposition. It gives only
$\|Q_m\|_{\tr}\le r+(m-1)/d$ for Schmidt rank $r$, which exceeds the tight value by $(r-1)(d+1-m)/d$
(Remark~\ref{rem:triangle}). The factorisation above avoids this loss.

Both main theorems have been formalised in Lean~4~\cite{Lean4} with Mathlib~\cite{Mathlib2020}
(Section~\ref{sec:lean}).

\section{Setting and results}\label{sec:results}

Throughout, $d\ge2$, vectors live in $\C^d$ with its standard basis, and $\bar v$ denotes the entrywise complex
conjugate of a vector $v$. The trace norm $\|A\|_{\tr}=\|A\|_1$ of a matrix is the sum of its singular values,
and $\|A\|_2=(\Tr A^\dagger A)^{1/2}$ is the Hilbert--Schmidt (Frobenius) norm. For a matrix $X$ we write
$|X|=(X^\dagger X)^{1/2}$ and $|X^\dagger|=(XX^\dagger)^{1/2}$.

\begin{definition}
\begin{enumerate}
\item \emph{MUBs.} A family $\{e^k_a\}$, $k=1,\dots,m$, $a=1,\dots,d$, consists of $m$ mutually unbiased bases if
each $\{e^k_a\}_a$ is an orthonormal basis of $\C^d$ and $|\langle e^k_a|e^l_b\rangle|^2=1/d$ for $k\neq l$.
\item \emph{Equiangular measurements (EAMs).} Unit vectors $\psi_1,\dots,\psi_n\in\C^d$ with $n>d$ form an
equiangular measurement if $\sum_a\ket{\psi_a}\bra{\psi_a}=\frac nd\one$ and
$|\langle\psi_a|\psi_b\rangle|^2=c$ for $a\neq b$, where necessarily $c=c_{n,d}=\frac{n-d}{d(n-1)}$. The
measurement operators are $E_a=\frac dn\ket{\psi_a}\bra{\psi_a}$.
\item \emph{Schmidt number.} A state $\rho$ on $\C^d\otimes\C^d$ has Schmidt number at most $r$, written
$\SN(\rho)\le r$, if it is a finite convex combination of projectors onto unit vectors of Schmidt rank at most
$r$~\cite{TerhalHorodecki2000}.
\item \emph{Correlation matrices}~\cite{TavakoliMorelli2024}. For MUBs $\{e^k_a\}$ (Alice) and $\{f^l_b\}$ (Bob),
$Q_m(\rho)$ is the real $md\times md$ matrix with entries
$Q_{(k,a),(l,b)}=\bra{e^k_a\otimes f^l_b}\rho\ket{e^k_a\otimes f^l_b}$. For EAMs $\{\psi_a\}$ and $\{\varphi_b\}$,
$P_n(\rho)$ is the $n\times n$ matrix with entries
$P_{ab}=\Tr[\rho(E_a\otimes F_b)]=\frac{d^2}{n^2}\bra{\psi_a\otimes\varphi_b}\rho\ket{\psi_a\otimes\varphi_b}$.
\end{enumerate}
\end{definition}

The value of $c_{n,d}$ follows by taking the trace of $(\sum_a\ket{\psi_a}\bra{\psi_a})^2=\frac{n^2}{d^2}\one$.

\begin{theorem}[Conjecture~3 of \cite{TavakoliMorelli2024}]\label{thm:mub}
Let $\{e^k_a\}$ and $\{f^l_b\}$ be two sets of $m\ge1$ mutually unbiased bases of $\C^d$, and let
$\SN(\rho)\le r$. Then
\[
\|Q_m(\rho)\|_{\tr}\;\le\;1+\frac{(m-1)r}{d}.
\]
For every $d$, every $m$ for which $m$ MUBs exist in $\C^d$, and every $1\le r\le d$, equality is attained.
\end{theorem}

\begin{theorem}[Conjecture~2 of \cite{TavakoliMorelli2024}]\label{thm:eam}
Let $\{\psi_a\}$ and $\{\varphi_b\}$ be equiangular measurements with $n>d$ elements in $\C^d$, and let
$\SN(\rho)\le r$. Then
\[
\|P_n(\rho)\|_{\tr}\;\le\;\frac{d(d-1)+r(n-d)}{n(n-1)}.
\]
Equality is attained for $r\in\{1,d-1,d\}$ for every EAM, and for every $r$ whenever the real span of the
projectors $\ket{\psi_a}\bra{\psi_a}$ contains a projector of rank $r$ (for instance for SICs, $n=d^2$).
\end{theorem}

Both theorems are special cases of the following inequality. For a finite family $a=\{a_i\}$ of vectors in $\C^d$
(not necessarily normalised) and a Hermitian matrix $A$, write
\[
I_a(A)=\sum_i\bra{a_i}A\ket{a_i}^2 .
\]
\begin{definition}\label{def:frame}
A family $a$ has \emph{frame constants} $(\alpha,\beta)$, with $\alpha\ge\beta\ge0$, if for every Hermitian $A$
\begin{equation}\label{eq:frame}
I_a(A)\;\le\;\beta\Tr(A^2)+(\alpha-\beta)\frac{(\Tr A)^2}{d}.
\end{equation}
\end{definition}

\begin{theorem}[general bound]\label{thm:general}
Let $a=\{a_i\}_{i\in I}$ and $b=\{b_j\}_{j\in J}$ be families in $\C^d$ with frame constants $(\alpha_a,\beta_a)$
and $(\alpha_b,\beta_b)$, and let $M(\rho)$ be the $|I|\times|J|$ matrix with entries
$M_{ij}=\bra{a_i\otimes b_j}\rho\ket{a_i\otimes b_j}$. If $\SN(\rho)\le r$, then
\[
\|M(\rho)\|_{\tr}\le\Bigl(\beta_a+(\alpha_a-\beta_a)\frac rd\Bigr)^{1/2}\Bigl(\beta_b+(\alpha_b-\beta_b)\frac rd\Bigr)^{1/2}.
\]
More precisely, for a unit vector $\ket\psi=\sum_{ij}X_{ij}\ket i\ket j$,
\begin{equation}\label{eq:pure}
\|M(\ket\psi\bra\psi)\|_{\tr}\le\Bigl(\beta_a+(\alpha_a-\beta_a)\frac{\|X\|_1^2}{d}\Bigr)^{1/2}
\Bigl(\beta_b+(\alpha_b-\beta_b)\frac{\|X\|_1^2}{d}\Bigr)^{1/2}.
\end{equation}
\end{theorem}

Lemmas~\ref{lem:frameMUB} and~\ref{lem:frameEAM} below give the frame constants
\[
\text{$m$ MUBs: }(\alpha,\beta)=(m,1);\qquad
\text{EAM vectors }\sqrt{d/n}\,\psi_a:\ (\alpha,\beta)=\Bigl(\frac dn,\ \frac{d(d-1)}{n(n-1)}\Bigr).
\]
With these values Theorem~\ref{thm:general} gives exactly the bounds of Theorems~\ref{thm:mub}
and~\ref{thm:eam}. For the EAM,
$\beta+(\alpha-\beta)\frac rd=\frac{d(d-1)}{n(n-1)}+\frac rn-\frac{(d-1)r}{n(n-1)}=\frac{d(d-1)+r(n-d)}{n(n-1)}$.

\section{Proof of the general bound}\label{sec:proof}

\subsection{Pure states and coefficient matrices}
A vector $\ket\psi=\sum_{ij}X_{ij}\ket i\ket j$ of $\C^d\otimes\C^d$ is determined by its coefficient matrix
$X\in\C^{d\times d}$. Then $\|\psi\|^2=\|X\|_2^2$, the Schmidt coefficients of $\psi$ are the singular values of
$X$, and the Schmidt rank of $\psi$ is $\operatorname{rank}X$. For $u,v\in\C^d$,
\[
\langle u\otimes v|\psi\rangle=\sum_{ij}\bar u_i\bar v_jX_{ij}=\bra uX\ket{\bar v},
\]
so $M_{ij}(\ket\psi\bra\psi)=|\bra{a_i}X\ket{\bar b_j}|^2$.

\subsection{Gram factorisation}

\begin{lemma}\label{lem:factor}
Let $X\in\C^{d\times d}$ and let $\{a_i\}_{i\in I}$, $\{b_j\}_{j\in J}$ be finite families in $\C^d$. The matrix
$M_{ij}=|\bra{a_i}X\ket{b_j}|^2$ satisfies
\[
\|M\|_1\;\le\;I_a(|X^\dagger|)^{1/2}\,I_b(|X|)^{1/2}.
\]
\end{lemma}

\begin{proof}
Let $X=\sum_s\sigma_s\ket{u_s}\bra{v_s}$ be a singular value decomposition with $\sigma_s>0$ and orthonormal
families $\{u_s\}$, $\{v_s\}$. Put
\[
L=\sum_s\sigma_s^{1/2}\ket{u_s}\bra{u_s},\qquad R=\sum_s\sigma_s^{1/2}\ket{v_s}\bra{v_s},\qquad
W=\sum_s\ket{u_s}\bra{v_s}.
\]
Then $LWR=X$, $L^2=|X^\dagger|$, $R^2=|X|$, and $W^\dagger W R=R$, because $W^\dagger W=\sum_s\ket{v_s}\bra{v_s}$
is the projector onto the range of $R$. Let $x_i=La_i$ and $y_j=WRb_j$. Since $L$ is Hermitian,
$\bra{a_i}X\ket{b_j}=\langle x_i|y_j\rangle$. With $\xi_i=x_i\otimes\bar x_i$ and $\eta_j=y_j\otimes\bar y_j$ we get
$\langle\xi_i|\eta_j\rangle=\langle x_i|y_j\rangle\overline{\langle x_i|y_j\rangle}=M_{ij}$, that is,
$M=\Xi^\dagger H$, where $\Xi$ and $H$ are the matrices with columns $\xi_i$ and $\eta_j$. H\"older's inequality for
Schatten norms, $\|\Xi^\dagger H\|_1\le\|\Xi\|_2\|H\|_2$~\cite{Bhatia1997}, gives the claim, since
\[
\|\Xi\|_2^2=\sum_i\|\xi_i\|^2=\sum_i\|x_i\|^4=\sum_i\bra{a_i}L^2\ket{a_i}^2=I_a(|X^\dagger|)
\]
and
\[
\|H\|_2^2=\sum_j\|y_j\|^4=\sum_j\bra{b_j}RW^\dagger WR\ket{b_j}^2=\sum_j\bra{b_j}R^2\ket{b_j}^2=I_b(|X|). \qedhere
\]
\end{proof}

\subsection{Proof of Theorem~\ref{thm:general}}
Let $\ket\psi$ be a unit vector with coefficient matrix $X$, so $\|X\|_2=1$. By Lemma~\ref{lem:factor}, applied to
the families $a$ and $\bar b=\{\bar b_j\}$,
\[
\|M(\ket\psi\bra\psi)\|_1\le I_a(|X^\dagger|)^{1/2}I_{\bar b}(|X|)^{1/2}.
\]
The family $\bar b$ has the same frame constants as $b$: for Hermitian $A$ we have
$I_{\bar b}(A)=\sum_j\bra{b_j}\bar A\ket{b_j}^2=I_b(\bar A)$, and $\bar A$ is Hermitian with the same trace and the
same Hilbert--Schmidt norm as $A$. The matrices $|X^\dagger|$ and $|X|$ are Hermitian, with
$\Tr|X^\dagger|=\Tr|X|=\|X\|_1$ and $\Tr|X^\dagger|^2=\Tr|X|^2=\|X\|_2^2=1$. Inserting them into~\eqref{eq:frame}
gives~\eqref{eq:pure}. If $\psi$ has Schmidt rank at most $r$, then $X$ has at most $r$ nonzero singular values,
and the Cauchy--Schwarz inequality gives $\|X\|_1^2\le r\|X\|_2^2=r$. Since $\alpha\ge\beta$, the right-hand side
of~\eqref{eq:pure} is nondecreasing in $\|X\|_1^2$, which proves the bound for pure states. Finally, if
$\rho=\sum_tp_t\ket{\psi_t}\bra{\psi_t}$ with $p_t\ge0$, $\sum_tp_t=1$ and each $\psi_t$ of Schmidt rank at most
$r$, then $M(\rho)=\sum_tp_tM(\ket{\psi_t}\bra{\psi_t})$, and the triangle inequality for $\|\cdot\|_1$ completes
the proof. \qed

\section{Frame constants of MUBs and of equiangular measurements}\label{sec:frames}

We work in the real Hilbert space $\Herm(d)$ of Hermitian matrices with $\langle A,B\rangle=\Tr(AB)$. For a unit
vector $v$, $\bra vA\ket v=\langle P_v,A\rangle$ with $P_v=\ket v\bra v$. Every Hermitian $A$ decomposes
orthogonally as $A=\frac{\Tr A}{d}\one+A_0$ with $\Tr A_0=0$, and
$\|A_0\|_2^2=\Tr(A^2)-(\Tr A)^2/d$.

\begin{lemma}[MUB frames]\label{lem:frameMUB}
The family of all vectors of $m$ mutually unbiased bases has frame constants $(m,1)$:
\[
\sum_{k=1}^m\sum_{a=1}^d\bra{e^k_a}A\ket{e^k_a}^2\;\le\;\Tr(A^2)+(m-1)\frac{(\Tr A)^2}{d}.
\]
Equality holds if and only if $A_0\in S_1\oplus\dots\oplus S_m$, with $S_k$ as in the proof.
\end{lemma}

\begin{proof}
Fix $k$ and let $P^k_a=\ket{e^k_a}\bra{e^k_a}$. The $P^k_a$, $a=1,\dots,d$, are orthonormal in $\Herm(d)$ and span
$V_k=\R\one\oplus S_k$, where $S_k=\{\sum_ac_aP^k_a:\sum_ac_a=0\}$ is orthogonal to $\one$. Hence
\[
\sum_a\langle P^k_a,A\rangle^2=\|\Pi_{V_k}A\|_2^2=\frac{(\Tr A)^2}{d}+\|\Pi_{S_k}A_0\|_2^2 ,
\]
where $\Pi$ denotes orthogonal projection. For $k\neq l$ the subspaces $S_k$ and $S_l$ are orthogonal:
\[
\Tr\Bigl[\Bigl(\sum_ac_aP^k_a\Bigr)\Bigl(\sum_bc'_bP^l_b\Bigr)\Bigr]=\sum_{a,b}c_ac'_b\,|\langle e^k_a|e^l_b\rangle|^2
=\frac1d\Bigl(\sum_ac_a\Bigr)\Bigl(\sum_bc'_b\Bigr)=0 .
\]
By Bessel's inequality, $\sum_k\|\Pi_{S_k}A_0\|_2^2\le\|A_0\|_2^2=\Tr(A^2)-(\Tr A)^2/d$, with equality if and only
if $A_0\in\bigoplus_kS_k$. Summing over $k$ gives $I(A)\le m(\Tr A)^2/d+\Tr(A^2)-(\Tr A)^2/d$.
\end{proof}

\begin{lemma}[EAM frames]\label{lem:frameEAM}
Let $\psi_1,\dots,\psi_n$ be an equiangular measurement in $\C^d$, $n>d$, and $c=c_{n,d}$. For every Hermitian $A$,
\[
\sum_{a=1}^n\bra{\psi_a}A\ket{\psi_a}^2\;\le\;\frac n{d^2}(\Tr A)^2+(1-c)\Bigl(\Tr(A^2)-\frac{(\Tr A)^2}{d}\Bigr),
\]
with equality if $A$ lies in the real span of the projectors $P_a=\ket{\psi_a}\bra{\psi_a}$. Consequently the
vectors $\sqrt{d/n}\,\psi_a$ have frame constants $\bigl(\frac dn,\frac{d(d-1)}{n(n-1)}\bigr)$.
\end{lemma}

\begin{proof}
Let $T:\R^n\to\Herm(d)$, $T(x)=\sum_ax_aP_a$, with adjoint $(T^*A)_a=\langle P_a,A\rangle$. Then
$\sum_a\bra{\psi_a}A\ket{\psi_a}^2=\|T^*A\|^2=\langle A,TT^*A\rangle$. The Gram matrix is
$T^*T=G=(1-c)\one_n+cJ_n$, where $J_n$ is the all-ones matrix. It has eigenvalue $1-c+cn=n/d$ on the all-ones
vector $\mathbf 1_n$ and eigenvalue $1-c$ on $\mathbf 1_n^\perp$. Since $TT^*T=TG$:
\begin{itemize}
\item $T\mathbf 1_n=\sum_aP_a=\frac nd\one$ is an eigenvector of $TT^*$ with eigenvalue $n/d$;
\item $TT^*$ acts as $1-c$ on $V=T(\mathbf 1_n^\perp)=\{\sum_ax_aP_a:\sum_ax_a=0\}$;
\item $TT^*$ vanishes on $(\operatorname{range}T)^\perp$.
\end{itemize}
Every element of $V$ is traceless, so $V\perp\one$ and $TT^*=\frac nd\Pi_{\R\one}+(1-c)\Pi_V$. Therefore
\[
\langle A,TT^*A\rangle=\frac nd\cdot\frac{(\Tr A)^2}{d}+(1-c)\|\Pi_VA_0\|_2^2
\le\frac n{d^2}(\Tr A)^2+(1-c)\|A_0\|_2^2 .
\]
If $A$ lies in $\operatorname{range}T=\R\one\oplus V$, then $A_0\in V$ and equality holds. For the rescaled
vectors, multiply by $(d/n)^2$. Using $1-c=\frac{n(d-1)}{d(n-1)}$,
\[
\frac{d^2}{n^2}\Bigl[\frac n{d^2}(\Tr A)^2+(1-c)\Bigl(\Tr A^2-\frac{(\Tr A)^2}{d}\Bigr)\Bigr]
=\beta\Tr(A^2)+(\alpha-\beta)\frac{(\Tr A)^2}{d},
\]
with $\alpha=d/n$ and $\beta=\frac{d(d-1)}{n(n-1)}$, and $\alpha\ge\beta$ because $n>d$.
\end{proof}

\begin{proof}[Proof of the bounds in Theorems~\ref{thm:mub} and~\ref{thm:eam}]
Apply Theorem~\ref{thm:general} with the frame constants of Lemmas~\ref{lem:frameMUB} and~\ref{lem:frameEAM}. For
EAMs, use $a_i=\sqrt{d/n}\,\psi_i$ and $b_j=\sqrt{d/n}\,\varphi_j$, so that $M(\rho)=P_n(\rho)$.
\end{proof}

\section{Tightness}\label{sec:tight}

\begin{proposition}\label{prop:tight}
\begin{enumerate}
\item Let $\{e^k_a\}$ be $m$ MUBs in $\C^d$ and $1\le r\le d$. Let Bob use the conjugate bases
$f^l_b=\bar e^l_b$ (which are again MUBs). Let $\Pi$ be the projector onto the span of $e^1_1,\dots,e^1_r$, and
let $\psi$ have coefficient matrix $X=\Pi/\sqrt r$. Then $\psi$ is a unit vector of Schmidt rank $r$ and
$\|Q_m(\ket\psi\bra\psi)\|_{\tr}=1+(m-1)r/d$.
\item Let $\{\psi_a\}$ be an EAM, let Bob use $\varphi_b=\bar\psi_b$, and let $\Pi$ be a rank-$r$ projector in the
real span of the $P_a$. The state with $X=\Pi/\sqrt r$ attains the bound of Theorem~\ref{thm:eam}. Rank-$r$
projectors in this span exist for $r=1$ ($P_a$), $r=d$ ($\one=\frac dn\sum_aP_a$) and $r=d-1$ ($\one-P_a$), and
for every $r$ when the $P_a$ span $\Herm(d)$, as for SICs.
\end{enumerate}
\end{proposition}

\begin{proof}
In both cases $M_{ij}=|\bra{a_i}X\ket{a_j}|^2$ with $X\succeq0$, where $a$ is Alice's family (rescaled by
$\sqrt{d/n}$ in case 2). With $x_i=X^{1/2}a_i$ we get $M_{ij}=\langle x_i\otimes\bar x_i|x_j\otimes\bar x_j\rangle$.
So $M$ is a Gram matrix, hence positive semidefinite, and $\|M\|_1=\Tr M=\sum_i\bra{a_i}X\ket{a_i}^2=I_a(X)$.

In case~1, $\bra{e^1_a}\Pi\ket{e^1_a}\in\{0,1\}$ with $r$ ones, and
$\bra{e^k_a}\Pi\ket{e^k_a}=\sum_{b\le r}|\langle e^k_a|e^1_b\rangle|^2=r/d$ for $k\ne1$. So
$I(\Pi/\sqrt r)=\frac1r\bigl(r+(m-1)d\frac{r^2}{d^2}\bigr)=1+\frac{(m-1)r}{d}$.

In case~2, equality holds in Lemma~\ref{lem:frameEAM} because $X$ lies in the span of the $P_a$. With
$\Tr X=\sqrt r$ and $\Tr X^2=1$, the value of $I_a(X)$ is the right-hand side of~\eqref{eq:frame} at
$(\Tr A)^2=r$, which is the bound.
\end{proof}

\begin{remark}[intermediate Schmidt numbers for equiangular measurements]
For the simplex EAM ($n=d+1$) the span of the $P_a$ has dimension $d+1$, and the construction above does not
reach $2\le r\le d-2$. Numerical optimisation over pure states of Schmidt rank $r$ gives maxima strictly below
the bound in all cases we tried: $0.69496<0.7$ for $(d,r)=(4,2)$; $0.73100<0.73333$ and $0.76026<0.76667$ for
$d=5$, $r=2,3$; and $0.78230<0.78571$ for $(d,r)=(6,3)$. For $r=d-1$ the bound is attained, as proved above. We do not claim tightness at intermediate $r$ for general EAMs.
\end{remark}

\section{Remarks}\label{sec:remarks}

\begin{remark}[a stronger inequality]
For a pure state with coefficient matrix $X$ (not necessarily normalised), the proof gives
$\|Q_m\|_{\tr}\le\|X\|_2^2+(m-1)\|X\|_1^2/d$. Theorem~\ref{thm:mub} follows from this via
$\|X\|_1^2\le r\|X\|_2^2$.
\end{remark}

\begin{remark}[unequal sets and dimensions]
The proof of Theorem~\ref{thm:general} applies verbatim to $\C^{d_A}\otimes\C^{d_B}$, with the frame
inequality~\eqref{eq:frame} in the respective dimension. For $m_A$ MUBs in $\C^{d_A}$ and $m_B$ MUBs in $\C^{d_B}$
it gives
$\|Q\|_{\tr}\le\bigl[(1+(m_A-1)r/d_A)(1+(m_B-1)r/d_B)\bigr]^{1/2}$ for $\SN(\rho)\le r$.
\end{remark}

\begin{remark}[comparison with the triangle inequality]\label{rem:triangle}
Write a unit vector of Schmidt rank $r$ as $\psi=\sum_s\sqrt{\lambda_s}\,u_s\otimes v_s$. Expanding
$Q_m(\ket\psi\bra\psi)$ in the Schmidt terms and using the triangle inequality leads to bounds of the form
$\|X\|_1^2+(m-1)\|X\|_2^2/d$, that is, $r+(m-1)/d$ in the worst case. This exceeds
$1+(m-1)r/d$ by $(r-1)(d+1-m)/d$, which is positive for $r\ge2$ and $m\le d$. The Gram factorisation of
Lemma~\ref{lem:factor} places the trace norm $\|X\|_1^2$ on the ``coherent'' term $(m-1)/d$ and the
Hilbert--Schmidt norm on the diagonal term, which is the tight arrangement.
\end{remark}

\begin{remark}[scope]
Conjecture~1 of~\cite{TavakoliMorelli2024}, on complete sets of MUBs and SICs, is not addressed here.
\end{remark}

\section{Formal verification}\label{sec:lean}

Theorems~\ref{thm:mub} and~\ref{thm:eam} are formalised in Lean~4 (version 4.33.1)~\cite{Lean4} with the Mathlib
library~\cite{Mathlib2020} (revision \texttt{0df444a3}). The formal statements are:
\begin{alltt}\footnotesize
theorem tavakoli_morelli_conjecture_3 (hm : 1 \(\le\) m) (e f : Fin m \(\to\) Fin d \(\to\) Fin d \(\to\) \(\mathbb{C}\))
    (he : IsMUB d m e) (hf : IsMUB d m f) (\(\rho\) : Matrix (Fin d \(\times\) Fin d) (Fin d \(\times\) Fin d) \(\mathbb{C}\))
    (h\(\rho\) : HasSchmidtNumberLE \(\rho\) r) :
    traceNorm (probMatrix e f \(\rho\)) \(\le\) 1 + ((m : \(\mathbb{R}\)) - 1) * r / d

theorem tavakoli_morelli_conjecture_2 (hdn : d < n) \{\(\psi\) \(\varphi\) : Fin n \(\to\) Fin d \(\to\) \(\mathbb{C}\)\}
    (h\(\psi\) : IsEAM d n \(\psi\)) (h\(\varphi\) : IsEAM d n \(\varphi\)) (\(\rho\) : Matrix (Fin d \(\times\) Fin d) (Fin d \(\times\) Fin d) \(\mathbb{C}\))
    (h\(\rho\) : HasSchmidtNumberLE \(\rho\) r) :
    traceNorm (eamProbMatrix \(\psi\) \(\varphi\) \(\rho\)) \(\le\) ((d : \(\mathbb{R}\)) * (d - 1) + r * (n - d)) / (n * (n - 1))
\end{alltt}
The definitions are as follows.
\begin{itemize}
\item \lean{traceNorm} is the sum of the singular values, defined through the eigenvalues of $Q^{\mathsf T}Q$.
\item \lean{HasSchmidtNumberLE} states that $\rho$ is a finite convex combination of projectors onto unit vectors
of Schmidt rank at most $r$.
\item \lean{IsMUB} and \lean{IsEAM} are the definitions of Section~\ref{sec:results}.
\item \lean{probMatrix} and \lean{eamProbMatrix} are the matrices $Q_m$ and $P_n$.
\end{itemize}
The formal proof proceeds through the variational characterisation of the trace norm. It shows, using a singular
value decomposition built from Mathlib's spectral theorem, that $\|Q\|_{\tr}\le\gamma$ as soon as
$\Tr(O^{\mathsf T}Q)\le\gamma$ for every real orthogonal $O$. The bound on $\Tr(O^{\mathsf T}Q)$ (theorems
\lean{tavakoli\_morelli\_conjecture\_3\_dual} and \lean{\_2\_dual}) is proved along the lines of
Lemma~\ref{lem:factor} and Lemmas~\ref{lem:frameMUB} and~\ref{lem:frameEAM}.

All main theorems depend only on Lean's standard axioms (\lean{propext}, \lean{Classical.choice},
\lean{Quot.sound}). The development contains no \lean{sorry} and no \lean{native\_decide}. As sanity checks the
repository also contains non-vacuity examples (qubit MUBs and a Bell state).

\section*{Data availability}
The Lean project, the numerical checks and the verification logs are available at
\url{https://github.com/anshM123/Tavakoli-Morelli-Conjecture}.

\begin{thebibliography}{99}

\bibitem{TavakoliMorelli2024}
A.~Tavakoli and S.~Morelli, \emph{Enhanced Schmidt number criteria based on correlation trace norms},
Phys. Rev. A \textbf{110}, 062417 (2024), arXiv:2402.09972.

\bibitem{TerhalHorodecki2000}
B.~M.~Terhal and P.~Horodecki, \emph{Schmidt number for density matrices}, Phys. Rev. A \textbf{61}, 040301(R)
(2000).

\bibitem{Spengler2012}
C.~Spengler, M.~Huber, S.~Brierley, T.~Adaktylos, and B.~C.~Hiesmayr, \emph{Entanglement detection via mutually
unbiased bases}, Phys. Rev. A \textbf{86}, 022311 (2012).

\bibitem{Bavaresco2018}
J.~Bavaresco, N.~Herrera Valencia, C.~Kl\"ockl, M.~Pivoluska, P.~Erker, N.~Friis, M.~Malik, and M.~Huber,
\emph{Measurements in two bases are sufficient for certifying high-dimensional entanglement}, Nat. Phys.
\textbf{14}, 1032 (2018).

\bibitem{WoottersFields1989}
W.~K.~Wootters and B.~D.~Fields, \emph{Optimal state-determination by mutually unbiased measurements}, Ann. Phys.
\textbf{191}, 363 (1989).

\bibitem{Durt2010}
T.~Durt, B.-G.~Englert, I.~Bengtsson, and K.~\.Zyczkowski, \emph{On mutually unbiased bases}, Int. J. Quantum Inf.
\textbf{8}, 535 (2010).

\bibitem{Bhatia1997}
R.~Bhatia, \emph{Matrix Analysis}, Graduate Texts in Mathematics \textbf{169} (Springer, New York, 1997).

\bibitem{Lean4}
L.~de~Moura and S.~Ullrich, \emph{The Lean 4 theorem prover and programming language}, in \emph{Automated
Deduction -- CADE 28}, Lecture Notes in Computer Science \textbf{12699} (Springer, 2021), pp.~625--635.

\bibitem{Mathlib2020}
The mathlib Community, \emph{The Lean mathematical library}, in \emph{Proceedings of the 9th ACM SIGPLAN
International Conference on Certified Programs and Proofs (CPP 2020)} (ACM, 2020), pp.~367--381.

\end{thebibliography}

\end{document}
